% ============================================================
\section{Analysis and Discussion}
\label{sec:analysis}
% ============================================================

\textbf{Why does the C$>$D gap collapse at scale?}
At 200/5, Gaussian-noise negatives (C) outperform hard negatives (D)
because ``easy'' off-manifold negatives provide unambiguous gradient signal
at limited scale.  At 1000/20 with NLL-only ES, the gap collapses to
near-zero ($p{=}0.39$, n.s.), consistent with easy negatives providing clearer gradient signal when
training budget is limited, but this advantage becoming unnecessary
with sufficient data.
The detailed mechanism (Appendix~\ref{app:gaussian_mechanism}) shows that
easy negatives degrade \emph{most} at scale ($-$0.043), while hard negatives
are under-trained at 200/5 and improve at 500/10 ($+$0.068 for D).

\textbf{Catastrophic forgetting of pre-trained features.}
The 500/10$\to$1000/20 degradation is explained by catastrophic forgetting:
extended fine-tuning overwrites pre-trained features, collapsing Moirai
to CNN-level performance.
Freezing the encoder recovers AUC (frozen B$=$0.606, C$=$0.582, D$=$0.573
vs.\ unfrozen 0.506--0.521), confirming the mechanism.
The mechanism is implicit regularisation loss: more synthetic training data drives the encoder further from its pre-trained initialisation, until the resulting features are no better than random initialisation (CNN AUC$=$0.513$\approx$unfrozen Moirai; full analysis in Appendix~\ref{app:forgetting_mechanism}).

\textbf{Does Moirai pre-training help?}
At 200/5, the from-scratch CNN (AUC$=$0.597) outperforms all Moirai
conditions.
At 1000/20, frozen Moirai substantially exceeds both unfrozen Moirai
and CNN (0.513): pre-trained features \emph{are} useful when preserved.
\emph{Methodological caveat}: the CNN uses MSE reconstruction loss while
Moirai uses distributional NLL; the CNN advantage may partly reflect MSE
being an easier optimisation target, not just the absence of harmful
pre-training.  A definitive comparison would require training the CNN with a distributional
NLL head (e.g., Gaussian NLL); this controlled experiment is future work
(Appendix~\ref{app:cnn_analysis}).
Moirai's primary contribution to this work is diagnostic: the frozen-encoder
experiment requires pre-trained features whose destruction can be measured.
A CNN has no pre-trained features to lose, so it cannot reveal the
catastrophic forgetting mechanism.
The backbone-agnostic claim rests on two architectures (Moirai, CNN) and
one cross-domain validation (N-BaIoT); testing a third backbone
(e.g., PatchTST as encoder) would further strengthen Contribution~2.

\textbf{Recommended operating point.}
Default to the 1D-CNN for IoT anomaly detection at small scale
($\leq$200 samples): it achieves the highest AUC (0.597), lowest variance,
and has no pre-training dependency, though the MSE-vs-NLL loss confound
means this advantage is preliminary.
When pre-trained features are available and training budget exceeds
200 samples, switch to frozen Moirai (freeze the encoder, train only
the projection head, use NLL-only early stopping, recalibrate on
deployment-specific benign traffic).
\emph{Caveat}: neither configuration achieves operationally useful F1 at
stealth-95 under calibrated thresholds (best: C F1$=$0.250,
frozen C F1$=$0.148); the system is a research prototype best suited as
a high-specificity alerting layer, not a standalone detector.
LoRA~\cite{hu2022lora} and adapters~\cite{houlsby2019adapters} offer
a principled middle ground between full freezing and full fine-tuning;
this is future work.
Frozen-encoder AUC recovery does not translate to F1 recovery under
calibrated thresholds (frozen C: F1$=$0.148 vs.\ 200/5 C: F1$=$0.250);
threshold recalibration is needed operationally
(Appendix~\ref{app:freezing_details}).

\textbf{When to prefer \ours{} vs Ensemble.}
At stealth-95, the Ensemble (F1$=$0.136, FPR$=$0.042) and condition~D
(F1$=$0.137, FPR$=$0.040) are effectively tied, while condition~C
(F1$=$0.250, FPR$=$0.080) outperforms both at higher FPR.
\ours{} is preferable for
(1)~\emph{novel attack generalisation} (leave-one-out: mean $+$0.186 F1;
Table~\ref{tab:loo}) and
(2)~\emph{adversarial augmentation} without real attack traffic.
The Ensemble's partial robustness reflects our evaluation construction: its Isolation Forest component detects distributional shifts in the sparse feature subsets perturbed by each archetype, but this signal weakens monotonically with stealth level; real multi-feature attacks may differ (Appendix~\ref{app:ensemble_degradation}).

\textbf{Computational overhead.}
Hard-negative generation takes $<$2s/sample on CPU.
Fine-tuning Moirai-Small for 5 epochs takes $\approx$2--4 min on CPU;
inference is $<$50ms/sample.

\textbf{Limitations.}
We identify seven limitations (full details in Appendix~\ref{app:limitations}):
(1)~hand-designed archetypes;
(2)~three training scales with 5 seeds;
(3)~fixed 12-dim input, stealth 85--95\%;
(4)~Diffusion-TS backbone deferred;
(5)~circular evaluation (addressed by N-BaIoT, \S\ref{sec:experiments:nbaiot});
(6)~low absolute F1 (best 0.250; system is a high-specificity alerting layer);
(7)~baselines use published defaults.

\textbf{Responsible disclosure.}
We release code and data for defensive research, not to enable attacks.
